\chapter{Hemoglobin A1c Test}

\section*{Overview}
The hemoglobin A1c test measures the percentage of hemoglobin with glucose attached and estimates average blood glucose over approximately two to three months.

\section*{Common Instruments and Equipment}
The exact setup varies by technique and setting. Common items may include:
\begin{itemize}[nosep,leftmargin=*]
  \item venipuncture supplies and labeled blood collection tube
  \item point-of-care or central laboratory analyzer
  \item quality-control materials for the assay
\end{itemize}

\section*{Why It Is Done}
It is commonly used to diagnose diabetes or prediabetes and to monitor glucose control in people with diabetes.

\section*{How It Is Performed}
A blood sample is collected from a vein or fingerstick. Fasting is usually unnecessary, although other tests collected at the same visit may require fasting.

\section*{Contraindications}
No single contraindication applies to every patient; the treating team assesses whether the procedure is suitable for the individual and the reason for treatment.

\section*{Precautions}
The care team reviews relevant health conditions, medicines, allergies, and procedure-specific risks, and explains any changes needed for the individual.

\section*{Results and Limitations}
Results are reported as a percentage or in mmol/mol. Anemia, hemoglobin variants, kidney disease, pregnancy, transfusion, and altered red-cell survival can affect interpretation.

\section*{References}
\begin{enumerate}
  \item MedlinePlus. Hemoglobin A1c (HbA1c) Test. U.S. National Library of Medicine; 2025.
  \item Current Medical Diagnosis and Treatment. McGraw-Hill; 2025.
\end{enumerate}
\clearpage
